\documentclass[11pt]{article}
\usepackage{fontspec}
\usepackage{xeCJK}
\usepackage{xurl}
\usepackage[a4paper,margin=25mm]{geometry}
\usepackage{amsmath,amssymb,amsthm,mathtools}
\usepackage{xcolor}
\usepackage{booktabs,array,enumitem,fancyhdr,hyperref}
\hypersetup{colorlinks=true,linkcolor=blue!50!black,urlcolor=blue!50!black}
\setlength{\emergencystretch}{2em}
\setlength{\parskip}{0.3em}
\linespread{1.16}
\allowdisplaybreaks
\newtheorem{theorem}{定理}[section]
\newtheorem{proposition}[theorem]{命题}
\newtheorem{lemma}[theorem]{引理}
\theoremstyle{definition}
\newtheorem{definition}[theorem]{定义}
\theoremstyle{remark}
\newtheorem{remark}[theorem]{说明}
\newcommand{\dd}{\,\mathrm d}
\newcommand{\ii}{\mathrm i}
\newcommand{\Rea}{\operatorname{Re}}
\newcommand{\Ima}{\operatorname{Im}}
\newcommand{\Li}{\operatorname{Li}}
\newcommand{\supp}{\operatorname{supp}}
\pagestyle{fancy}
\fancyhf{}
\fancyhead[L]{\small 学习笔记}
\fancyhead[R]{\small 初稿}
\fancyfoot[C]{\thepage}
\setlength{\headheight}{15pt}
\author{}
\date{2026年10月22日初稿}
\makeatletter
\renewcommand{\maketitle}{%
  \thispagestyle{plain}%
  \begin{center}
    {\LARGE\bfseries \@title\par}
    \vspace{0.6em}
    {\small \@date\par}
  \end{center}
  \vspace{0.6em}}
\makeatother
\newcommand{\draftnotice}{\noindent{\small
本文为供逐篇审阅的学习笔记初稿，不是新增研究成果。
文中区分直接证明与引用结果；正文不依赖本文。
本轮已作公式、条件及排版初步核查，尚待作者逐句审定。}\par\medskip}

\XeTeXlinebreaklocale "zh"
\renewcommand{\proofname}{证明}
\newcommand{\Ei}{\operatorname{Ei}}
\newcommand{\li}{\operatorname{li}}
\date{2026年10月24日初稿}
\title{Fourier 部分和与零点截断：\\相同项数为何未必同样准确}
\begin{document}
\maketitle

\noindent{\small 本书第 18 章将阶梯函数的周期 Fourier 部分和与零点项截断并列，并指出两图“不能据此直接比较逼近效率”。
本篇把比较放在精确的设定下：对象取 Riemann 的 $J_0$（其显式公式没有模型误差），
先给出数值事实，再把差别分解为可以分别论证的几个因素。
凡标为命题者给出证明；引用结果注明出处；数值结果单独标明。}\par\medskip

\section{设定}
\subsection{两种逼近}
记
\[
 J_0(x)=\sum_{p^k<x}\frac1k+\frac12\sum_{p^k=x}\frac1k.
\]
在非跳跃点第二个和为空，在素数幂处这一定义取左右平均。
本篇用 $\li(x)=\Ei(\log x)$ 表示通常对数积分，不是从 $2$ 起积分的偏移函数 $\Li(x)$。
Riemann 显式公式为：对 $x>1$，
\[
 J_0(x)=\li(x)-\sum_\rho\Ei(\rho\log x)-\log2+\int_x^\infty\frac{\dd t}{t(t^2-1)\log t},       \tag{1.1}
\]
其中 $\Ei$ 采用本书第 14 章的分支约定，零点按高度对称截断、共轭成对求和。
(1.1) 是恒等式，因此下面的截断误差完全来自截断本身。

\begin{itemize}
 \item \textbf{零点逼近}：取前 $K$ 个正虚部零点 $\rho_1,\ldots,\rho_K$，令
 \[
  J^{(K)}(x)=\li(x)-2\sum_{k=1}^K\Rea\Ei(\rho_k\log x)-\log2+\int_x^\infty\frac{\dd t}{t(t^2-1)\log t}.
 \]
 \item \textbf{Fourier 逼近}：将 $J_0$ 在 $[a,b)$ 上的限制作周期 $L=b-a$ 的延拓，取 $N$ 阶部分和
 \[
  S_N(x)=c_0+\sum_{n=1}^N\Bigl[a_n\cos\frac{2\pi n(x-a)}L+b_n\sin\frac{2\pi n(x-a)}L\Bigr].
 \]
\end{itemize}
“相同项数”在本文只约定为 $K=N$：比较 $K$ 对共轭零点与 $N$ 阶 Fourier 部分和。
一对共轭零点在实公式中合并为一个 $2\Rea\Ei$ 项，
这一定义不意味着两种方案具有相同自由度、相同计算成本或相同输入信息。
Fourier 系数预先使用了已知的阶梯函数，因此本实验比较表示效果，不比较独立求素数的算法效率。

\subsection{数值事实}
在观察窗口 $[13,35]$ 上，以步长 $0.05$ 取 $441$ 个点，计算均方根误差 $\mathrm{rms}$，
以及距跳跃点（$13,16,17,19,23,25,27,29,31,32$）不小于 $0.5$ 的点上的最大误差 $\max'$。
所用零点与 $\Ei$ 由 \texttt{mpmath} 计算，Fourier 系数由分段常数函数精确积分得到。
所用前400个正虚部零点位于本书第19章所述的已认证高度范围内，
使用它们的临界线位置不等于假设全部零点满足 RH；本程序的坐标和函数求值仍为普通数值近似。

\begin{center}
\small
\begin{tabular}{lll@{\hspace{2.5em}}lll}
\toprule
零点逼近 & rms & $\max'$ & Fourier 逼近 & rms & $\max'$\\
\midrule
$K=25$  & $0.136$ & $0.168$ & $[10,40)$, $N=25$  & $0.137$ & $0.267$\\
$K=50$  & $0.104$ & $0.106$ & $[10,40)$, $N=100$ & $0.040$ & $0.041$\\
$K=100$ & $0.072$ & $0.076$ & $[2,100)$, $N=25$  & $0.289$ & $0.536$\\
$K=200$ & $0.047$ & $0.048$ & $[2,100)$, $N=100$ & $0.107$ & $0.175$\\
$K=250$ & $0.039$ & $0.043$ & $[2,1000)$, $N=25$  & $10.19$ & $15.76$\\
$K=400$ & $0.026$ & $0.030$ & $[2,1000)$, $N=100$ & $3.06$ & $5.64$\\
\bottomrule
\end{tabular}
\end{center}

由表可见：
\begin{enumerate}
 \item 在第 18 章的设定（周期区间 $[10,40)$，$N=100$）下，Fourier 逼近确实更好；
 零点逼近要达到同样的 rms，约需 $K=250$。
 \item 把 Fourier 的周期区间放大到 $[2,100)$，$N=100$ 的误差已劣于 $K=100$ 的零点逼近；
 放大到 $[2,1000)$ 时，本次固定 $N$ 的误差进一步增大。
\end{enumerate}
因此“Fourier 更优”不是两类函数系的固有性质，而依赖于比较的设定。
本书第18章的两图使用 $\pi$ 与简化的零点余弦模型；
本篇改用同一个精确对象 $J_0$，是另设的受控比较，不是对原两图误差的直接定量归因。

\section{截断核与局部振荡尺度}
\subsection{Fourier：误差由跳跃与频宽决定}
\begin{proposition}\label{prop:fourier}
设 $f$ 是周期 $L$ 的分段常数函数，在一个周期内的跳跃点为 $t_j$，跳跃量为 $h_j$。则在非跳跃点 $x$ 处
\[
 |S_Nf(x)-f(x)|\le\sum_j\frac{|h_j|}{\pi(N+1)\,|\sin(\pi(x-t_j)/L)|}.
\]
\end{proposition}
\begin{proof}
令 $s(t)=\{t/L\}-\tfrac12$，它在 $L\mathbb Z$ 处跳跃 $-1$，斜率为 $1/L$。由于 $\sum_jh_j=0$（周期性），
$f(x)=c-\sum_jh_js(x-t_j)$，其中 $c$ 为常数。由 $s(t)=-\sum_{n\ge1}\sin(2\pi nt/L)/(\pi n)$，
\[
 S_Nf(x)-f(x)=-\sum_jh_j\sum_{n>N}\frac{\sin(2\pi n(x-t_j)/L)}{\pi n}.
\]
对 $\theta\notin2\pi\mathbb Z$，$|\sum_{n=M}^{M'}\sin n\theta|\le1/|\sin(\theta/2)|$，而 $1/n$ 单调下降，
由分部求和（学习笔记第 7 篇）得 $|\sum_{n>N}\sin(n\theta)/n|\le1/((N+1)|\sin(\theta/2)|)$。
取 $\theta=2\pi(x-t_j)/L$ 即得。
\end{proof}
在离跳跃点 $t_j$ 距离为 $d\ll L$ 处，单项的界与 $|h_j|/(\Omega_Fd)$ 同阶（相差因子 $2/\pi$），其中
\[
 \Omega_F=\frac{2\pi N}L
\]
是 $S_N$ 所含的最高角频率（对变量 $x$）。误差由“跳跃量 / (频宽 $\times$ 距离)”控制。
周期延拓在 $a\equiv b$ 处还有一个跳跃 $J_0(a)-J_0(b^-)$，它同样计入上式。

\subsection{零点：与截断 Perron 积分的联系}
令 $a_n=\Lambda(n)/\log n$，对 $\Rea s>1$ 有 $\log\zeta(s)=\sum a_nn^{-s}$。
\begin{lemma}[截断 Perron 公式，引用]\label{lem:perron}
对 $y>0$，$y\ne1$，$c>0$，$T>0$，
\[
 \Bigl|\frac1{2\pi\ii}\int_{c-\ii T}^{c+\ii T}\frac{y^s}s\dd s-\mathbf 1_{y>1}\Bigr|<y^c\min\Bigl(1,\frac1{T|\log y|}\Bigr).
\]
（Davenport, \emph{Multiplicative Number Theory}, Ch.~17。）
\end{lemma}
\begin{proposition}\label{prop:perron}
设 $c>1$，$x>1$ 不是素数幂。令 $I_T(x)=\frac1{2\pi\ii}\int_{c-\ii T}^{c+\ii T}\log\zeta(s)\frac{x^s}s\dd s$。则
\[
 |I_T(x)-J_0(x)|\le\sum_{n\ge2}a_n\Bigl(\frac xn\Bigr)^c\min\Bigl(1,\frac1{T|\log(x/n)|}\Bigr).
\]
\end{proposition}
\begin{proof}
在 $\Rea s=c$ 上 $\sum a_nn^{-c}<\infty$，可以逐项积分；对每一项用引理~\ref{lem:perron}，取 $y=x/n$。
\end{proof}
$I_T$ 沿 $s=c+\ii t$ 截取 $|t|\le T$；在对数变量 $u=\log x$ 中，
去掉因子 $e^{cu}$ 后可视为有限频带的 Fourier 积分。
命题~\ref{prop:perron} 的单项界还带权 $a_n(x/n)^c$，其中与距离有关的因子是
$1/(T|u-\log n|)$，与命题~\ref{prop:fourier} 的结构相似。
把积分线左移并处理零点处的对数割线，可得到主项与高度内零点的 $\Ei$ 项，水平线段另有误差；
对 $\psi$ 的完整论证见 Davenport 第 17 章，$J$ 的情形结构相同（零点处为对数支点，相应留数换为 $\Ei$ 项）。
本篇没有定量估计 $J^{(K)}-I_T$，故上述不等式首先是 $I_T$ 的误差界，
不能直接改写为 $J^{(K)}$ 的误差界。
零点项含有缓变幅度，也不构成严格的有限带宽 Fourier 级数。

回到变量 $x$：在跳跃点 $n$ 附近，$|\log(x/n)|\approx|x-n|/x$，命题~\ref{prop:perron} 的单项界约为 $1/((T/x)|x-n|)$。
从相位 $\gamma\log x$ 的导数看，零点截断在 $x$ 处的局部角频率尺度为
\[
 \Omega_Z(x)=\frac{\gamma_K}x .
\]

\subsection{频宽的比较}
$\gamma_{100}=236.524\ldots$，$L=30$。于是
\[
 \Omega_F=\frac{2\pi\cdot100}{30}=20.94,\qquad
 \Omega_Z(13)=18.19,\quad\Omega_Z(24)=9.86,\quad\Omega_Z(35)=6.76 .
\]
在整个窗口内 $\Omega_F>\Omega_Z(x)$，在窗口中部约为两倍。
这一差别为表中的误差提供了有用的解释线索：同样的 $K=N=100$，
在该窗口内 Fourier 的最高角频率大于零点项的局部相位变化尺度。
但上界不能当作实际误差渐近，且 $J^{(K)}-I_T$ 尚未估计，
所以这里不据频率比推出误差必相差两倍，也不认定已证明某因素占主要作用。

\section{第二个因素：局部表示与全局表示}
Fourier 部分和只表示 $J_0$ 在一个长为 $L$ 的周期区间上的限制；
零点截断采用同一公式定义于全部 $x>1$，但这不保证固定 $K$ 时在整个半轴上误差一致小。
两者的频率分配方式不同：
\begin{itemize}
 \item Fourier 的频宽 $2\pi N/L$ 在区间内均匀，区间越长，同样 $N$ 的频宽越低；
 \item 零点的频宽 $\gamma_K/x$ 随 $x$ 反比下降，在对数尺度 $u=\log x$ 中均匀。
\end{itemize}
这解释了表中 Fourier 一栏随区间增大而迅速恶化：$[2,1000)$ 上 $N=100$ 的频宽只有 $0.63$，低于零点逼近在 $x\le375$ 处的频宽。

零点的频宽与项数的关系由零点计数决定。由 $N(T)=\frac T{2\pi}\log\frac T{2\pi e}+O(\log T)$（引用），
$K$ 个零点达到的频宽 $\gamma_K$ 满足 $K\approx\frac{\gamma_K}{2\pi}\log\frac{\gamma_K}{2\pi e}$。
平均而言，零点在高处越来越密，所以高度覆盖随项数的增长比线性略慢；
它与等距 Fourier 频率的指标增长方式不同。

若进一步比较分辨跳跃所需的项数，还须规定误差范数与跳点排除邻域。
平均素数间距不能替代最小素数幂间距，有限个连续函数之和也不能在跳跃两侧任意一致逼近阶梯函数。
因此本文不从上述频率尺度推出两种方案分辨全部跳跃所需项数的普遍阶数。

\section{第三个因素：系数是否局部最优}
\begin{proposition}\label{prop:proj}
在 $L^2(a,b)$ 中，$S_Nf$ 是 $f$ 在 $\{1,\cos\frac{2\pi n(x-a)}L,\sin\frac{2\pi n(x-a)}L:\ n\le N\}$ 张成的子空间上的正交投影，
因而是该子空间中 $L^2$ 误差最小的元素。
\end{proposition}
\begin{proof}
这些函数两两正交；Fourier 系数恰为正交投影的系数。由勾股定理，子空间中任一元素 $g$ 满足
$\|f-g\|^2=\|f-S_Nf\|^2+\|S_Nf-g\|^2$。
\end{proof}
上述最优性属于整个周期区间 $(a,b)$ 的连续 $L^2$ 范数，
不是子窗口 $[13,35]$ 上离散采样误差的最优性。
零点逼近的系数由 (1.1) 全局确定，并非对窗口最优：函数 $\Rea\Ei(\rho_k\log x)$、$\Ima\Ei(\rho_k\log x)$ 在有限区间上一般不正交，
截断也不是投影。可以取 $1$、式(1.1)的固定主项与修正项之和，
以及 $2K$ 个 $\Rea$、$\Ima\Ei(\rho_k\log x)$，重新拟合系数。
但这是另一个允许自由系数的模型，不再是原显式公式的截断。
这些列在有限窗口可能近似线性相关，拟合结果须同时报告数值秩、奇异值截断与精度稳定性。
配套程序输出不同截断阈值下的秩和残差，作为数值诊断；
本文不把一次双精度最小二乘的改善比例当作排除某个误差因素的证明。

\section{第四个因素：远处跳跃的贡献}
命题~\ref{prop:fourier} 中只出现一个周期内的跳跃。命题~\ref{prop:perron} 则对一切素数幂 $n$ 求和，
并带有权 $(x/n)^c$：小素数幂距离 $x$ 远，单项界中的 $1/(T|\log(x/n)|)$ 较小，但权 $(x/n)^c$ 较大。
这表明截断 Perron 积分的上界涉及整个算术序列，而周期 Fourier 界涉及一个周期内的跳跃，
包括周期拼接处的跳跃。上界中的各项不构成实际误差的可加分解，
尤其在未估计 $I_T-J^{(K)}$ 时，不能给零点截断的实测误差分配这些项的百分比。

\Needspace{18\baselineskip}
\section{结论}
\begin{enumerate}
 \item 两种截断的核与相位提供了比较分辨尺度的方法；本文已证明 Fourier 尾和界与截断 Perron 积分界，
 尚未将后者定量转换为零点截断误差界。
 \item 本次 $L=30$、$K=N=100$ 实验中 Fourier 误差较小；其最高角频率较高
 （$20.9$ 对零点项局部尺度 $6.8$--$18.2$），是合理的解释线索而非误差比例定理。
 \item 这一优势依赖于窗口：Fourier 只表示一个周期区间，零点截断的同一公式定义于全部 $x>1$，
 其相位频率以对数变量计量。在本次所列扩大周期区间的实验中，原有优势消失。
 \item 正交投影、自由系数拟合与远处跳跃均涉及不同的比较对象；本篇不据现有实验给它们确定普遍主次关系。
 \item 因此，两种展开的差别不在于“效率高低”，而在于它们回答不同的问题：
Fourier 部分和是在给定周期区间的相应三角多项式空间中的 $L^2$ 最优表示；
显式公式则是由零点联系全体素数幂的全局恒等式。
\end{enumerate}

\section*{参考与数据}
\begin{itemize}
 \item H.\ Davenport, \emph{Multiplicative Number Theory}, 3rd ed., Springer, 2000, Ch.~17.
 \item H.\ M.\ Edwards, \emph{Riemann's Zeta Function}, Academic Press, 1974, Ch.~1.
 \item H.\ Riesel, G.\ Göhl, Some calculations related to Riemann's prime number formula, \emph{Math.\ Comp.} 24 (1970).
 \item 本书第 14 章（Riemann 显式公式）、第 18 章（零点项的对数尺度振荡）；学习笔记第 7 篇。
\end{itemize}
表中数据由本目录的 \texttt{计算比较数据.py} 生成（\texttt{mpmath} 30 位工作精度，零点由 \texttt{zetazero} 计算），
是数值结果，不是区间认证。程序在临时缓存目录保存不同精度的数据，
30位与50位工作精度下，表中舍入数值一致；计算 $\Ei$ 后转为双精度汇总，
这一比较只核查所列显示精度，不是任意位数或整个区间的严格误差保证。
最小二乘结果单独标为诊断，不以提高 $\Ei$ 工作精度代替对线性代数病态性的分析。
\end{document}
